% Zone „Das kennst du schon“ – aus bank/funktionsklassen-und-eigenschaften/zone.jsonl (zusammenbau v0.1)
\einheitenkopf[zone][Kennst du schon]{Das kennst du schon}
\setcounter{aufgabe}{0}

%% TODO Titel: Ich-kann-Satz fehlt in der Bank (Platzhalter: Kettenname) – Nr. 1
%% TODO Anweisung: ein Satz über der ersten Teilaufgabe fehlt in der Bank – Nr. 1
\begin{aufgabe}{Termwerte berechnen}
\begin{teile}
\teil $f(x) = x^2 + 4$ – wie groß ist $f(3)$? f(3) = \leerfeld
\teil $h(x) = \frac{1}{2}x^2 - 3x$ – wie groß ist $h(-4)$? h(−4) = \leerfeld
\end{teile}
\end{aufgabe}

%% TODO Titel: Ich-kann-Satz fehlt in der Bank (Platzhalter: Kettenname) – Nr. 2
%% TODO Anweisung: ein Satz über der ersten Teilaufgabe fehlt in der Bank – Nr. 2
\begin{aufgabe}{Lineare Funktionen}
\begin{teile}
\teil $y = 3x - 6$ – Nullstelle? x = \leerfeld
\teil Gerade durch $A(-2 | 1)$ und $B(2 | 3)$ – Steigung? m = \leerfeld
\end{teile}
\end{aufgabe}

%% TODO Titel: Ich-kann-Satz fehlt in der Bank (Platzhalter: Kettenname) – Nr. 3
%% TODO Anweisung: ein Satz über der ersten Teilaufgabe fehlt in der Bank – Nr. 3
\begin{aufgabe}{Maßstab und Einheiten}
\begin{teile}
\teil 1 Längeneinheit entspricht 2 m. 6 Längeneinheiten – wie viele Meter? \leerfeld[m]
\teil 1 Längeneinheit entspricht 0,5 m. 3,4 Längeneinheiten – wie viele Zentimeter? \leerfeld[cm]
\end{teile}
\end{aufgabe}

%% TODO Titel: Ich-kann-Satz fehlt in der Bank (Platzhalter: Kettenname) – Nr. 4
%% TODO Anweisung: ein Satz über der ersten Teilaufgabe fehlt in der Bank – Nr. 4
\begin{aufgabe}{Quadratische Gleichungen lösen}
\begin{gleichungsraster}[2]
\gl{\text{$x^2 = 49$ – Lösungen?}} &
\gl{\text{$x^2 - 2x - 8 = 0$ – Lösungen?}} \\
\end{gleichungsraster}
\end{aufgabe}

%% TODO Titel: Ich-kann-Satz fehlt in der Bank (Platzhalter: Kettenname) – Nr. 5
%% TODO Anweisung: ein Satz über der ersten Teilaufgabe fehlt in der Bank – Nr. 5
\begin{aufgabe}{Die Grundgraphen kennen}
\begin{teile}
\teil Wie ist die Parabel $y = -x^2$ geöffnet? \\ \kreuz{nach oben} \\ \kreuz{nach unten}
\teil Welche Werte nimmt $y = \mathrm{sin}(x)$ an? \\ \kreuz{alle reellen Zahlen} \\ \kreuz{alle Zahlen von $-1$ bis $1$} \\ \kreuz{nur Zahlen größer als null}
\end{teile}
\end{aufgabe}

%% TODO Titel: Ich-kann-Satz fehlt in der Bank (Platzhalter: Kettenname) – Nr. 6
%% TODO Anweisung: ein Satz über der ersten Teilaufgabe fehlt in der Bank – Nr. 6
\begin{aufgabe}{Spiegeln im Koordinatensystem}
\begin{teile}
\teil $P(3 | 5)$ an der x-Achse gespiegelt – Bildpunkt? P'( \leerfeld | \leerfeld )
\teil $R(1{,}5 | -3)$ an der y-Achse gespiegelt – Bildpunkt? R'( \leerfeld | \leerfeld )
\end{teile}
\end{aufgabe}

%% TODO Titel: Ich-kann-Satz fehlt in der Bank (Platzhalter: Kettenname) – Nr. 7
%% TODO Anweisung: ein Satz über der ersten Teilaufgabe fehlt in der Bank – Nr. 7
\begin{aufgabe}{Termwerte berechnen – Fehler finden}
\begin{teile}
\teil Ich finde den Fehler: $r(x) = x^3 - 6x$ an der Stelle $-3$. \rechnung{r(-3) &= (-3)^3 - 6 \cdot (-3) \\ &= 27 + 18 \\ &= 45}
\end{teile}
\end{aufgabe}

%% TODO Titel: Ich-kann-Satz fehlt in der Bank (Platzhalter: Kettenname) – Nr. 8
%% TODO Anweisung: ein Satz über der ersten Teilaufgabe fehlt in der Bank – Nr. 8
\begin{aufgabe}{Termwerte berechnen – selbst rechnen}
\begin{teile}
\teil $s(x) = x^3 - 5x$ – wie groß ist $s(-2)$? s(−2) = \leerfeld
\end{teile}
\end{aufgabe}
